\chapter{Getting started}
\label{ch:start}

\section{Before you start}

You need three things:

\begin{enumerate}
  \item \textbf{An account on \server.} If you don't have one yet, ask a bot admin. You get a
        personal signup link; it shows your username (taken from your university email
        address) and asks for a password.
  \item \textbf{Element.} Use the desktop app, the phone app or the web version. On the login
        screen, change the homeserver to \server{} (\ui{Edit} next to the server name) and
        log in with your username and password.
  \item \textbf{The professor role.} A bot admin adds you with \cmd{!admin prof add}. TAs
        don't need this: the course owner appoints them.
\end{enumerate}

\begin{warning}[Keep two things: your recovery key and your reset code]
\begin{itemize}
  \item \textbf{Recovery key} (from Element). When Element asks you to set up
        \emph{recovery} or \emph{secure backup}, do it and keep the key in a password
        manager or on paper. Course rooms and your chat with the bot are end-to-end
        encrypted. Without the key, a new device or a password reset means you can't read
        old messages any more, and nobody can get them back for you.
  \item \textbf{Reset code} (from the bot). The signup page shows it once. With it you can
        set a new password yourself if you forget yours, even when you are logged out
        everywhere. Lost it? Send \cmd{!password code} to the bot for a new one.
\end{itemize}
Keep both \emph{outside} Matrix: when you need them, you can't read your chats.
\autoref{sec:passwords} explains what to do if you lose one of them.
\end{warning}

\section{Start a chat with the bot}

\begin{enumerate}
  \item In Element, click \ui{+} next to \ui{People} (or \ui{Start new chat}).
  \item Enter \botid{} and start the chat.
  \item The bot joins within a few seconds. Send \cmd{!whoami}.
\end{enumerate}

\begin{you}
!whoami
\end{you}
\begin{botreply}
You are @prof.example:matrix.dmi.unibas.ch.\\
Roles: professor\\
Courses: -\\[4pt]
\cmd{!help} lists the commands these roles allow.
\end{botreply}

If the bot only offers \cmd{!password} and \cmd{!password code}, you don't have a role yet:
for everyone who isn't a professor, TA or admin, the bot only handles their own password
(\autoref{sec:no-reply}).

\begin{note}[Only in your direct chat]
The bot accepts commands only in a chat that you and the bot are alone in. Commands can
contain email addresses and signup links, which must not end up in a shared room. If you
type a command in a course room, the bot asks you to use the direct chat instead. Don't
invite other people into your chat with the bot.
\end{note}

\section{How the bot answers}
\label{sec:answers}

\begin{description}[style=nextline,leftmargin=1.5em,font=\normalfont\bfseries]
  \item[\robot\ The robot reaction]
    As soon as the bot starts working on your command, it reacts to your message with
    \robot. The reply follows when the work is done, usually within a few seconds.
    No reaction means the bot hasn't received the command yet (\autoref{sec:slow}).
  \item[Plans and \cmd{!confirm}]
    Commands that change many things at once, such as inviting 200 students, don't act
    right away. The bot first replies with a \emph{plan}: what it would do, and a short
    code. Nothing has happened at that point. You run the plan with \cmd{!confirm <code>} or
    drop it with \cmd{!cancel}. Plans expire after 10~minutes.
  \item[Jobs]
    A confirmed plan runs as a \emph{job} in the background. The bot posts one progress
    message, updates it, and sends a summary at the end, including anything that failed
    and what to do next. You can keep using other commands meanwhile.
  \item[Files]
    Some commands need a file: an email list, a CSV table, a template or a bot plugin.
    Send the command first; the bot then asks for the file and waits 10~minutes for it.
    Alternatively, upload the file first and reply to it with the command. If the bot
    rejects a file, it says why and waits for a corrected one: you only send the file
    again, not the command.
\end{description}

\begin{you}
!invite CS101 --dry-run
anna.beispiel@stud.unibas.ch
\end{you}
\begin{botreply}
Invite to \textbf{CS101}: 0 existing accounts get an invitation, 1 new people get a signup
link (as a CSV file after you confirm), 0 entries are not usable.\\[4pt]
\emph{Dry run: nothing was changed. Send the same command without \cmd{--dry-run} to get a
plan you can confirm.}
\end{botreply}

\section{Getting help in the chat}

\begin{itemize}
  \item \cmd{!help}: all commands you can use, grouped by topic.
  \item \cmd{!help course}: the commands of one group (here: courses) with a short description.
  \item \cmd{!help course create}: one command in detail, with arguments, an example, the
        next steps and who may use it.
  \item \cmd{!course help} and \cmd{!course create --help} do the same.
\end{itemize}

When you mistype a command, the bot suggests the closest one (\enquote{Did you mean
\cmd{!course create}?}). When arguments are missing, it shows the correct form and points to
the detailed help.

\section{Your first course in six steps}
\label{sec:quickstart}

\begin{enumerate}
  \item \textbf{Create the course.}
\begin{you}
!course create CS101 "Introduction to Computer Science" --semester HS26
\end{you}
        Accept the invitations to the new space and its rooms in Element.
  \item \textbf{Appoint your TAs} (they need an account already):
\begin{you}
!course ta add CS101 @ta.one:matrix.dmi.unibas.ch
\end{you}
  \item \textbf{Invite the students.} Paste the email addresses below the command, or send
        the list from the course platform as a file:
\begin{you}
!invite CS101
anna.beispiel@stud.unibas.ch
max.muster@stud.unibas.ch
\end{you}
        Check the plan and send \cmd{!confirm <code>}.
  \item \textbf{Send the signup links.} The bot sends you \file{signup-links.csv} with one
        link per new person. Send them by mail merge (\autoref{sec:mailmerge}).
  \item \textbf{Create exercise groups} when the groups are known:
\begin{you}
!groups create CS101 8
\end{you}
        and after \cmd{!confirm}: \cmd{!groups assign CS101 --auto}.
  \item \textbf{At the end of the semester}: \cmd{!course rollover CS101 FS27} for the next
        run and \cmd{!course archive CS101-HS26} for the old one (\autoref{sec:lifecycle}).
\end{enumerate}
